% Generated draft: requires mathematical and source review.
\documentclass[11pt]{article}
\usepackage[margin=25mm]{geometry}
\usepackage{fontspec}
\setmainfont{Latin Modern Roman}
\usepackage{amsmath,amssymb,mathtools,unicode-math}
\setmathfont{Latin Modern Math}
\usepackage{enumitem,needspace,xurl,hyperref}
\setlength{\parindent}{0pt}
\setlength{\parskip}{5pt}
\setlength{\emergencystretch}{3em}
\newcommand{\N}{\mathbb{N}}
\newcommand{\Z}{\mathbb{Z}}
\newcommand{\Q}{\mathbb{Q}}
\newcommand{\R}{\mathbb{R}}
\newcommand{\C}{\mathbb{C}}
\newcommand{\F}{\mathbb{F}}
\newcommand{\A}{\mathbb{A}}
\newcommand{\PP}{\mathbb{P}}
\newcommand{\E}{\mathbb{E}}
\newcommand{\eps}{\varepsilon}
\newcommand{\dd}{\,\mathrm{d}}
\begin{document}
\section*{UnsolvedMath ID: 11000191}
\textbf{Source problem code:} catalogue entry
\subsection{Definitions and mathematical statement}
Let \(\rho:\pi_1(\Sigma)\to G\) be a representation, and let \([\rho]\) denote its class in the source's representation/character space \(X\). The mapping class action is induced by precomposition with surface automorphisms (up to inner automorphisms), for example \(f\cdot[\rho]=[\rho\circ f_*^{-1}]\). The task is to give necessary and sufficient conditions on \(\rho\) for
\[\overline{\operatorname{Mod}_\Sigma\cdot[\rho]}^{\,X}=X.\]
Thus it seeks a characterization of individual dense orbits, not merely ergodicity or almost-everywhere density. Density of \(\rho(\pi_1\Sigma)\) in \(G\) is another property and is not asserted equivalent to orbit density.

The ambient \(X\), topology, group class and surface conventions are defined only in omitted context. For compact \(G\), a natural model is the topological conjugacy quotient of \(\operatorname{Hom}(\pi_1\Sigma,G)\); for other groups a chosen character quotient and component may be essential. If \(\Sigma\) has boundary and mapping classes preserve the boundary labels, the conjugacy classes of boundary holonomies are invariant. Consequently one must specify whether \(X\) is the relative space with those fixed boundary values, a component or some other invariant locus; orbit density in that locus and density in the full character space are different questions. No boundary data or connected-component target is silently chosen.

The source expressly singles out \(G=SU(2)\) and an \(n\)-holed sphere \(\Sigma_{0,n}\) with \(n>4\), saying that case remains open. This is preserved as a historical source qualification, not independently asserted as current status. Boundary marking/permutation, allowed holonomies and relative-space conventions are missing even in that stated special case.
\subsection{Short English statement}
Characterize exactly which representation classes have dense mapping-class orbits in the intended character-space locus. The source highlights \(SU(2)\) on spheres with more than four holes, but omits the ambient topology and boundary conditions.
\subsection{Sources}
\begin{itemize}
\item Local UnsolvedMath registry, permanent ID \(11000191\), exact source code \texttt{catalogue entry}.
\item Farb, \emph{Problems on Mapping Class Groups and Related Topics}, Problem 2.4, PDF page 218: \url{https://www.math.uchicago.edu/~farb/papers/mcgbook.pdf}. Registry access 2026-07-29; pdf-font-structure extraction.
\item The research summary invokes Goldman's mapping-class/character-variety theory without a precise bibliographic or resolving reference.
\end{itemize}
\end{document}
